% ============================================================
% fig_dependencias.tex — Dependencias largas y correferencia
% Uso: % ============================================================
% fig_dependencias.tex — Dependencias largas y correferencia
% Uso: % ============================================================
% fig_dependencias.tex — Dependencias largas y correferencia
% Uso: % ============================================================
% fig_dependencias.tex — Dependencias largas y correferencia
% Uso: \input{figures/fig_dependencias}
% ============================================================
\begin{figure}[H]
\centering
\begin{tikzpicture}[
    word/.style={draw=gray!35, rounded corners, fill=gray!8,
                 minimum width=1.05cm, minimum height=0.62cm,
                 font=\scriptsize, inner sep=2pt},
    subj/.style={draw=blue!60, rounded corners, fill=blue!12,
                 minimum width=1.05cm, minimum height=0.62cm,
                 font=\scriptsize\bfseries, inner sep=2pt},
    verb/.style={draw=red!60, rounded corners, fill=red!10,
                 minimum width=1.05cm, minimum height=0.62cm,
                 font=\scriptsize\bfseries, inner sep=2pt},
    pron/.style={draw=orange!70, rounded corners, fill=orange!12,
                 minimum width=1.05cm, minimum height=0.62cm,
                 font=\scriptsize\bfseries, inner sep=2pt},
    lbl/.style={font=\tiny, align=center}
]

% ── Tokens colocados manualmente (12 tokens) ─────────────────
\node[subj] (t0)  at (0*1.08, 3.2) {El};
\node[subj] (t1)  at (1*1.08, 3.2) {cient.};
\node[word] (t2)  at (2*1.08, 3.2) {que};
\node[word] (t3)  at (3*1.08, 3.2) {gan\'o};
\node[word] (t4)  at (4*1.08, 3.2) {el};
\node[word] (t5)  at (5*1.08, 3.2) {premio};
\node[word] (t6)  at (6*1.08, 3.2) {Nobel};
\node[word] (t7)  at (7*1.08, 3.2) {de};
\node[word] (t8)  at (8*1.08, 3.2) {f\'isica};
\node[verb] (t9)  at (9*1.08, 3.2) {present\'o};
\node[pron] (t10) at (10*1.08,3.2) {su};
\node[word] (t11) at (11*1.08,3.2) {invest.};

% ── Arco: sujeto → verbo (dependencia larga) ─────────────────
\draw[-{Stealth}, blue!60, thick]
    (t1.north) .. controls (1.08, 4.6) and (9.72, 4.6) .. (t9.north);
\node[lbl, blue!70] at (5.4, 4.85)
    {dependencia sujeto--verbo \\ (8 tokens de distancia)};

% ── Arco: pronombre → sujeto (correferencia) ─────────────────
\draw[-{Stealth}, orange!80, thick]
    (t10.south) .. controls (10.8, 1.8) and (1.08, 1.8) .. (t1.south);
\node[lbl, orange!70] at (5.4, 1.4)
    {\textit{``su''} refiere a \textit{``cient.''} \\ (correferencia a distancia)};

% ── Leyenda ──────────────────────────────────────────────────
\node[subj, label=right:{\scriptsize sujeto}]    at (0.0,  0.55) {};
\node[verb, label=right:{\scriptsize verbo}]     at (2.5,  0.55) {};
\node[pron, label=right:{\scriptsize pronombre}] at (5.0,  0.55) {};

% ── Nota ─────────────────────────────────────────────────────
\node[draw=gray!30, fill=gray!5, rounded corners, font=\tiny,
      inner sep=5pt, align=left] at (9.5, 0.55)
    {Un MLP con ventana fija no puede\\ conectar tokens tan alejados.\\ Un Transformer s\'i: la atenci\'on es global.};

\end{tikzpicture}
\caption{Dependencia de largo alcance y correferencia. El verbo
\textit{present\'o} concuerda con \textit{cient\'ifico}, su sujeto,
separados por 8 tokens. El pronombre \textit{su} refiere tambi\'en al mismo
sujeto. Ambas relaciones requieren acceder a contexto distante.}
\label{fig:dependencias}
\end{figure}

% ============================================================
\begin{figure}[H]
\centering
\begin{tikzpicture}[
    word/.style={draw=gray!35, rounded corners, fill=gray!8,
                 minimum width=1.05cm, minimum height=0.62cm,
                 font=\scriptsize, inner sep=2pt},
    subj/.style={draw=blue!60, rounded corners, fill=blue!12,
                 minimum width=1.05cm, minimum height=0.62cm,
                 font=\scriptsize\bfseries, inner sep=2pt},
    verb/.style={draw=red!60, rounded corners, fill=red!10,
                 minimum width=1.05cm, minimum height=0.62cm,
                 font=\scriptsize\bfseries, inner sep=2pt},
    pron/.style={draw=orange!70, rounded corners, fill=orange!12,
                 minimum width=1.05cm, minimum height=0.62cm,
                 font=\scriptsize\bfseries, inner sep=2pt},
    lbl/.style={font=\tiny, align=center}
]

% ── Tokens colocados manualmente (12 tokens) ─────────────────
\node[subj] (t0)  at (0*1.08, 3.2) {El};
\node[subj] (t1)  at (1*1.08, 3.2) {cient.};
\node[word] (t2)  at (2*1.08, 3.2) {que};
\node[word] (t3)  at (3*1.08, 3.2) {gan\'o};
\node[word] (t4)  at (4*1.08, 3.2) {el};
\node[word] (t5)  at (5*1.08, 3.2) {premio};
\node[word] (t6)  at (6*1.08, 3.2) {Nobel};
\node[word] (t7)  at (7*1.08, 3.2) {de};
\node[word] (t8)  at (8*1.08, 3.2) {f\'isica};
\node[verb] (t9)  at (9*1.08, 3.2) {present\'o};
\node[pron] (t10) at (10*1.08,3.2) {su};
\node[word] (t11) at (11*1.08,3.2) {invest.};

% ── Arco: sujeto → verbo (dependencia larga) ─────────────────
\draw[-{Stealth}, blue!60, thick]
    (t1.north) .. controls (1.08, 4.6) and (9.72, 4.6) .. (t9.north);
\node[lbl, blue!70] at (5.4, 4.85)
    {dependencia sujeto--verbo \\ (8 tokens de distancia)};

% ── Arco: pronombre → sujeto (correferencia) ─────────────────
\draw[-{Stealth}, orange!80, thick]
    (t10.south) .. controls (10.8, 1.8) and (1.08, 1.8) .. (t1.south);
\node[lbl, orange!70] at (5.4, 1.4)
    {\textit{``su''} refiere a \textit{``cient.''} \\ (correferencia a distancia)};

% ── Leyenda ──────────────────────────────────────────────────
\node[subj, label=right:{\scriptsize sujeto}]    at (0.0,  0.55) {};
\node[verb, label=right:{\scriptsize verbo}]     at (2.5,  0.55) {};
\node[pron, label=right:{\scriptsize pronombre}] at (5.0,  0.55) {};

% ── Nota ─────────────────────────────────────────────────────
\node[draw=gray!30, fill=gray!5, rounded corners, font=\tiny,
      inner sep=5pt, align=left] at (9.5, 0.55)
    {Un MLP con ventana fija no puede\\ conectar tokens tan alejados.\\ Un Transformer s\'i: la atenci\'on es global.};

\end{tikzpicture}
\caption{Dependencia de largo alcance y correferencia. El verbo
\textit{present\'o} concuerda con \textit{cient\'ifico}, su sujeto,
separados por 8 tokens. El pronombre \textit{su} refiere tambi\'en al mismo
sujeto. Ambas relaciones requieren acceder a contexto distante.}
\label{fig:dependencias}
\end{figure}

% ============================================================
\begin{figure}[H]
\centering
\begin{tikzpicture}[
    word/.style={draw=gray!35, rounded corners, fill=gray!8,
                 minimum width=1.05cm, minimum height=0.62cm,
                 font=\scriptsize, inner sep=2pt},
    subj/.style={draw=blue!60, rounded corners, fill=blue!12,
                 minimum width=1.05cm, minimum height=0.62cm,
                 font=\scriptsize\bfseries, inner sep=2pt},
    verb/.style={draw=red!60, rounded corners, fill=red!10,
                 minimum width=1.05cm, minimum height=0.62cm,
                 font=\scriptsize\bfseries, inner sep=2pt},
    pron/.style={draw=orange!70, rounded corners, fill=orange!12,
                 minimum width=1.05cm, minimum height=0.62cm,
                 font=\scriptsize\bfseries, inner sep=2pt},
    lbl/.style={font=\tiny, align=center}
]

% ── Tokens colocados manualmente (12 tokens) ─────────────────
\node[subj] (t0)  at (0*1.08, 3.2) {El};
\node[subj] (t1)  at (1*1.08, 3.2) {cient.};
\node[word] (t2)  at (2*1.08, 3.2) {que};
\node[word] (t3)  at (3*1.08, 3.2) {gan\'o};
\node[word] (t4)  at (4*1.08, 3.2) {el};
\node[word] (t5)  at (5*1.08, 3.2) {premio};
\node[word] (t6)  at (6*1.08, 3.2) {Nobel};
\node[word] (t7)  at (7*1.08, 3.2) {de};
\node[word] (t8)  at (8*1.08, 3.2) {f\'isica};
\node[verb] (t9)  at (9*1.08, 3.2) {present\'o};
\node[pron] (t10) at (10*1.08,3.2) {su};
\node[word] (t11) at (11*1.08,3.2) {invest.};

% ── Arco: sujeto → verbo (dependencia larga) ─────────────────
\draw[-{Stealth}, blue!60, thick]
    (t1.north) .. controls (1.08, 4.6) and (9.72, 4.6) .. (t9.north);
\node[lbl, blue!70] at (5.4, 4.85)
    {dependencia sujeto--verbo \\ (8 tokens de distancia)};

% ── Arco: pronombre → sujeto (correferencia) ─────────────────
\draw[-{Stealth}, orange!80, thick]
    (t10.south) .. controls (10.8, 1.8) and (1.08, 1.8) .. (t1.south);
\node[lbl, orange!70] at (5.4, 1.4)
    {\textit{``su''} refiere a \textit{``cient.''} \\ (correferencia a distancia)};

% ── Leyenda ──────────────────────────────────────────────────
\node[subj, label=right:{\scriptsize sujeto}]    at (0.0,  0.55) {};
\node[verb, label=right:{\scriptsize verbo}]     at (2.5,  0.55) {};
\node[pron, label=right:{\scriptsize pronombre}] at (5.0,  0.55) {};

% ── Nota ─────────────────────────────────────────────────────
\node[draw=gray!30, fill=gray!5, rounded corners, font=\tiny,
      inner sep=5pt, align=left] at (9.5, 0.55)
    {Un MLP con ventana fija no puede\\ conectar tokens tan alejados.\\ Un Transformer s\'i: la atenci\'on es global.};

\end{tikzpicture}
\caption{Dependencia de largo alcance y correferencia. El verbo
\textit{present\'o} concuerda con \textit{cient\'ifico}, su sujeto,
separados por 8 tokens. El pronombre \textit{su} refiere tambi\'en al mismo
sujeto. Ambas relaciones requieren acceder a contexto distante.}
\label{fig:dependencias}
\end{figure}

% ============================================================
\begin{figure}[H]
\centering
\begin{tikzpicture}[
    word/.style={draw=gray!35, rounded corners, fill=gray!8,
                 minimum width=1.05cm, minimum height=0.62cm,
                 font=\scriptsize, inner sep=2pt},
    subj/.style={draw=blue!60, rounded corners, fill=blue!12,
                 minimum width=1.05cm, minimum height=0.62cm,
                 font=\scriptsize\bfseries, inner sep=2pt},
    verb/.style={draw=red!60, rounded corners, fill=red!10,
                 minimum width=1.05cm, minimum height=0.62cm,
                 font=\scriptsize\bfseries, inner sep=2pt},
    pron/.style={draw=orange!70, rounded corners, fill=orange!12,
                 minimum width=1.05cm, minimum height=0.62cm,
                 font=\scriptsize\bfseries, inner sep=2pt},
    lbl/.style={font=\tiny, align=center}
]

% ── Tokens colocados manualmente (12 tokens) ─────────────────
\node[subj] (t0)  at (0*1.08, 3.2) {El};
\node[subj] (t1)  at (1*1.08, 3.2) {cient.};
\node[word] (t2)  at (2*1.08, 3.2) {que};
\node[word] (t3)  at (3*1.08, 3.2) {gan\'o};
\node[word] (t4)  at (4*1.08, 3.2) {el};
\node[word] (t5)  at (5*1.08, 3.2) {premio};
\node[word] (t6)  at (6*1.08, 3.2) {Nobel};
\node[word] (t7)  at (7*1.08, 3.2) {de};
\node[word] (t8)  at (8*1.08, 3.2) {f\'isica};
\node[verb] (t9)  at (9*1.08, 3.2) {present\'o};
\node[pron] (t10) at (10*1.08,3.2) {su};
\node[word] (t11) at (11*1.08,3.2) {invest.};

% ── Arco: sujeto → verbo (dependencia larga) ─────────────────
\draw[-{Stealth}, blue!60, thick]
    (t1.north) .. controls (1.08, 4.6) and (9.72, 4.6) .. (t9.north);
\node[lbl, blue!70] at (5.4, 4.85)
    {dependencia sujeto--verbo \\ (8 tokens de distancia)};

% ── Arco: pronombre → sujeto (correferencia) ─────────────────
\draw[-{Stealth}, orange!80, thick]
    (t10.south) .. controls (10.8, 1.8) and (1.08, 1.8) .. (t1.south);
\node[lbl, orange!70] at (5.4, 1.4)
    {\textit{``su''} refiere a \textit{``cient.''} \\ (correferencia a distancia)};

% ── Leyenda ──────────────────────────────────────────────────
\node[subj, label=right:{\scriptsize sujeto}]    at (0.0,  0.55) {};
\node[verb, label=right:{\scriptsize verbo}]     at (2.5,  0.55) {};
\node[pron, label=right:{\scriptsize pronombre}] at (5.0,  0.55) {};

% ── Nota ─────────────────────────────────────────────────────
\node[draw=gray!30, fill=gray!5, rounded corners, font=\tiny,
      inner sep=5pt, align=left] at (9.5, 0.55)
    {Un MLP con ventana fija no puede\\ conectar tokens tan alejados.\\ Un Transformer s\'i: la atenci\'on es global.};

\end{tikzpicture}
\caption{Dependencia de largo alcance y correferencia. El verbo
\textit{present\'o} concuerda con \textit{cient\'ifico}, su sujeto,
separados por 8 tokens. El pronombre \textit{su} refiere tambi\'en al mismo
sujeto. Ambas relaciones requieren acceder a contexto distante.}
\label{fig:dependencias}
\end{figure}
